En réponse à nos questions, notre \robot\ peut indiquer dans quelle direction il est orienté.
N'étant pas doué de la parole, ses réponses sont {\em booléennes} \index{Booléen} c'est-à-dire 
{\mcl True} (prononcer ``trou'', {\em vrai} \index{Vrai} en anglais) ou {\mcl False} (prononcer ``faulse'', 
{\em faux} \index{Faux} en anglais), écrits en \python\ avec une majuscule. On peut imaginer qu'il donne ses réponses à l'aide d'un voyant lumineux
qu'il allume pour répondre  {\mcl True} (vrai) et qu'il laisse éteint pour répondre {\mcl False}
(faux).

Pour poser les questions sur l'orientation du \robot, on utilise les commandes de test suivantes~:

\begin{table}[htbp]
\begin{center}
\begin{tabular}{@{}|l|p{11.7cm}|@{}} \hline
\multicolumn{2}{|l|}{\em Commandes de test d'orientation du \robot~:} \\ \hline \hline
{\idt tn()} & {\mcl True} si  le \robot\ est \underline{t}ourné vers le nord (\underline{n}) et {\mcl False} sinon. \\ \hline
{\idt tne()} & {\mcl True} si  le \robot\ est \underline{t}ourné vers le nord-est (\underline{ne}) et {\mcl False} sinon. \\ \hline
{\idt te()} & {\mcl True} si  le \robot\ est \underline{t}ourné vers l'est (\underline{e}) et {\mcl False} sinon. \\ \hline
{\idt tse()} & {\mcl True} si  le \robot\ est \underline{t}ourné vers le sud-est (\underline{se}) et {\mcl False} sinon. \\ \hline
{\idt ts()} & {\mcl True} si  le \robot\ est \underline{t}ourné vers le sud (\underline{s}) et {\mcl False} sinon. \\ \hline
{\idt tsw()} & {\mcl True} si  le \robot\ est \underline{t}ourné vers le sud-ouest (\underline{sw}) et {\mcl False} sinon. \\ \hline
{\idt tw()} & {\mcl True} si  le \robot\ est \underline{t}ourné vers l' ouest (\underline{w}) et {\mcl False} sinon. \\ \hline
{\idt tnw()} & {\mcl True} si  le \robot\ est \underline{t}ourné vers le nord-ouest (\underline{nw}) et {\mcl False} sinon. \\ \hline
\end{tabular}
\end{center}
%\label{tab:tnetc}
\caption{Tests d'orientation du \protect\robot}
\end{table}

En \python, l'instruction de test ``if'' permet de n'exécuter une séquence d'instructions que
si une certaine {\em expression booléenne} \index{Expression booléenne}est vraie ({\mcl True}).  Pour commencer nous allons
utiliser des expressions booléennes simples, à savoir les commandes de test d'orientation du \robot\
({\idt tn()}, {\idt tne()}, {\idt te()}, {\idt tse()}, {\idt ts()}, {\idt tsw()}, {\idt tw()}, {\idt tnw()}). Dans sa
forme la plus simple, l'instruction de test ``if'' s'écrit comme suit~:

\index{Test}
\begin{center}
\begin{tabular}{@{}|p{12.6cm}|@{}} \hline
{\em Test ``if'' à une seule alternative~:}\\ \hline \hline
\begin{minipage}[t]{12.6cm}
{\medskip
\mbox{\blc\ \ \ {\mcl if}\hspace{0.3em}\blanc\hspace{0.3em}{\small\rm expression booléenne}%
\footnote{\ A résultat {\mcl True} (vrai) ou {\mcl False} (faux).}{\idt :}}\par 
\mbox{\blc\ \ \ \ \ \ \ \small\rm  \ldots\ Instructions (indentées) exécutées seulement si}\par 
\mbox{\blc\ \ \ \ \ \ \ \small\rm  l'expression booléenne vaut {\mcl True} \ldots}\par 
\medskip}
\end{minipage}\\ \hline 
\end{tabular}
\end{center}

On peut traduire en français par {\idt si} \ldots\ {\idt alors}.
Le mot {\mcl if} se prononce `if' (si en anglais).
Comme pour le corps d'une procédure ou d'une boucle {\mcl for}, la ligne du {\mcl if} se termine par
un deux-points `{\idt :}' et les instructions qui dépendent du test sont indentées (décalées de quatre
blancs vers la droite par rapport au {\mcl if})~: elles forment un {\em bloc} qui se termine à la première ligne
moins indentée. Si l'on oublie le deux-points, \python\ refuse le programme avec le message
`{\tt SyntaxError: expected ':'}' (deux-points attendu) et si l'on oublie d'indenter,
avec le message
`{\tt IndentationError: expected an indented block after 'if' statement on line 5}'
(un bloc indenté était attendu après l'instruction {\mcl if} de la ligne 5, si le {\mcl if} est à la ligne 5 du programme).

Quand il n'y a qu'une seule instruction à exécuter, on peut l'écrire sur la même ligne, après le
deux-points, comme dans `{\mcl if}\ {\idt tw():\ pqtd()}'. Mais attention, dans
`{\mcl if}\ {\idt tw():\ pqtd();\ av()}' les deux instructions {\idt pqtd()} et {\idt av()} dépendent du test.
\underline{Une erreur grave} est de se tromper d'indentation. Par exemple dans~:
{
\parindent \parindentprog
\medskip

\vspace{\neglskip}\mbox{\blc {\mcl if}{\idt \ tw():}}\par
\vspace{\neglskip}\mbox{\blc \blk\blk{\idt pqtd()}}\par
\vspace{\neglskip}\mbox{\blc {\idt av()}}\par

\medskip }
\noindent seule l'instruction {\idt pqtd()} dépend du test~: l'instruction {\idt av()}
qui n'est pas indentée est toujours exé\-cu\-tée, que le \robot\ soit tourné vers l'ouest ou non.
\python\ ne peut pas signaler cette erreur puisque le programme est correct (mais ne fait pas
ce que l'on voulait). Notons aussi que l'instruction de test ``if'' est {\em une} instruction~:
si d'autres instructions la suivent, elles s'écrivent sur les lignes suivantes, avec la même
indentation que le {\mcl if}.


Dans l'exemple ci-dessous, le \robot\ avance deux fois et tourne à droite d'un quart de tour puis
d'un quart de tour supplémentaire 
s'il est orienté à l'ouest~:

\vspace{\baselineskip}

 {\raggedright 
 \includevspacePostScript[0.9cm]{FriseEnPeigne}{0.95}{0.95}
 \input python/FriseEnPeigne.tex
}
 
 Dans sa première forme ci-dessus l'instruction de test ne fait rien quand l'expression booléenne
 vaut {\mcl False}. Si l'on veut donner le choix de la séquence d'instructions à exécuter dans les
 deux alternatives, on utilisera la forme suivante~:
 
\index{Test}
 \begin{center}
\begin{tabular}{@{}|p{12.6cm}|@{}} \hline
{\em Test ``if'' à deux alternatives~:}\\ \hline \hline
\begin{minipage}[t]{12.6cm}
{\medskip
\mbox{\blc\ \ \ {\mcl if}\hspace{0.3em}\blanc\hspace{0.3em}{\small\rm expression booléenne}%
\footnote{ \ A résultat {\mcl True} (vrai) ou {\mcl False} (faux).}{\idt :}}\par 
\mbox{\blc\ \ \ \ \ \ \ \small\rm  \ldots\ Instructions (indentées) exécutées si l'expression}\par 
\mbox{\blc\ \ \ \ \ \ \ \small\rm   booléenne vaut {\mcl True} \ldots}\par 
\mbox{\blc\ \ \ {\mcl else}{\idt :}}\par 
\mbox{\blc\ \ \ \ \ \ \ \small\rm  \ldots\ Instructions  (indentées) exécutées si l'expression  }\par 
\mbox{\blc\ \ \ \ \ \ \ \small\rm  booléenne vaut {\mcl False} \ldots}\par 
\medskip}
\end{minipage}\\ \hline 
\end{tabular}
\end{center}

On peut traduire en français par {\idt si} \ldots\ {\idt alors} \ldots\ {\idt sinon} \ldots\
({\mcl else} se prononce `elsse' et signifie `autrement' en anglais).

Le {\mcl else} s'écrit exactement sous le {\mcl if} (avec la même indentation) et il est suivi
d'un deux-points. Comme précédemment, une alternative qui ne comporte qu'une seule instruction
peut s'écrire sur la même ligne que le {\mcl if} ou le {\mcl else}, après le deux-points.
C'est une erreur que d'oublier le deux-points après le {\mcl else} ou de mal aligner le
{\mcl else} sous le {\mcl if}, mais elle n'est pas
très grave car elle est signalée par \python\ avant l'exécution (`{\tt SyntaxError: expected ':'}'
ou `{\tt IndentationError: unindent does not match any outer indentation level}', le retrait
ne correspond à aucun niveau d'indentation englobant).
\python\ signale aussi une erreur (`{\tt SyntaxError: invalid syntax}') quand on oublie
d'indenter l'une des instructions placées entre le {\mcl if} et le {\mcl else}.

Dans le dessin de la frise en dents de scie ci-dessous, le \robot\ tourne de trois quarts de tour à droite 
s'il est orienté au nord-est et à gauche dans les autres cas (le seul autre cas possible étant d'ailleurs
qu'il soit tourné vers le sud)~:

\vspace{0.4cm}
 {\raggedright 
 \includevspacePostScript[0.5cm]{FriseEnDentsDeScie}{0.95}{0.95}
 \input python/FriseEnDentsDeScie.tex
 }
 \vfill
 Quand il y a plus de deux cas à tester on peut utiliser plusieurs alternatives et ne rien faire
 quand aucune des expressions booléeennes ne vaut {\mcl True} comme dans~:
 
\index{Test}
 \begin{center}
\begin{tabular}{@{}|p{12.6cm}|@{}} \hline
{\em Test ``if'' à $n$ alternatives~:}\\ \hline \hline
\begin{minipage}[t]{12.6cm}
{\medskip
\vspace{\neglskip}\mbox{\blc\ \ \ {\mcl if}\hspace{0.3em}\blanc\hspace{0.3em}{\small\rm expression booléenne 1}{\idt :}}\par 
\vspace{\neglskip}\mbox{\blc\ \ \ \ \ \ \ \small\rm  \ldots\ Instructions (indentées) exécutées si l'expression }\par 
\vspace{\neglskip}\mbox{\blc\ \ \ \ \ \ \ \small\rm  booléenne 1 vaut {\mcl True} \ldots}\par 
\vspace{\neglskip}\mbox{\blc\ \ \ {\mcl elif}\hspace{0.3em}\blanc\hspace{0.3em}{\small\rm expression booléenne 2}{\idt :}}\par 
\vspace{\neglskip}\mbox{\blc\ \ \ \ \ \ \ \small\rm  \ldots\ Instructions (indentées) exécutées si l'expression }\par 
\vspace{\neglskip}\mbox{\blc\ \ \ \ \ \ \ \small\rm  booléenne 2 vaut {\mcl True} \ldots}\par 
\mbox{\blc\ \ \ \dots}\par
\vspace{\neglskip}\mbox{\blc\ \ \ {\mcl elif}\hspace{0.3em}\blanc\hspace{0.3em}{\small\rm expression booléenne $n$}{\idt :}}\par 
\vspace{\neglskip}\mbox{\blc\ \ \ \ \ \ \ \small\rm  \ldots\ Instructions (indentées) exécutées si l'expression }\par 
\vspace{\neglskip}\mbox{\blc\ \ \ \ \ \ \ \small\rm  booléenne $n$ vaut {\mcl True} \ldots}\par 
\medskip}
\end{minipage}\\ \hline 
\end{tabular}
\end{center}

\vfill

 \noindent On peut aussi prévoir une dernière alternative qui est exécutée quand toutes les
 expressions booléeennes testées valent {\mcl False}~:
 
 \index{Test}
\begin{center}
\begin{tabular}{@{}|p{12.6cm}|@{}} \hline
{\em Test ``if'' à $(n + 1)$ alternatives~:}\\ \hline \hline
\begin{minipage}[t]{12.6cm}
{\medskip
\vspace{\neglskip}\mbox{\blc\ \ \ {\mcl if}\hspace{0.3em}\blanc\hspace{0.3em}{\small\rm expression booléenne 1}{\idt :}}\par 
\vspace{\neglskip}\mbox{\blc\ \ \ \ \ \ \ \small\rm  \ldots\ Instructions (indentées) exécutées si l'expression }\par 
\vspace{\neglskip}\mbox{\blc\ \ \ \ \ \ \ \small\rm  booléenne 1 vaut {\mcl True} \ldots}\par 
\vspace{\neglskip}\mbox{\blc\ \ \ {\mcl elif}\hspace{0.3em}\blanc\hspace{0.3em}{\small\rm expression booléenne 2}{\idt :}}\par 
\vspace{\neglskip}\mbox{\blc\ \ \ \ \ \ \ \small\rm  \ldots\ Instructions (indentées) exécutées si l'expression }\par 
\vspace{\neglskip}\mbox{\blc\ \ \ \ \ \ \ \small\rm  booléenne 2 vaut {\mcl True} \ldots}\par 
\mbox{\blc\ \ \ \dots}\par
\vspace{\neglskip}\mbox{\blc\ \ \ {\mcl elif}\hspace{0.3em}\blanc\hspace{0.3em}{\small\rm expression booléenne $n$}{\idt :}}\par 
\vspace{\neglskip}\mbox{\blc\ \ \ \ \ \ \ \small\rm  \ldots\ Instructions (indentées) exécutées si l'expression }\par 
\vspace{\neglskip}\mbox{\blc\ \ \ \ \ \ \ \small\rm  booléenne $n$ vaut {\mcl True} \ldots}\par 
\vspace{\neglskip}\mbox{\blc\ \ \ {\mcl else}{\idt :}}\par 
\vspace{\neglskip}\mbox{\blc\ \ \ \ \ \ \ \small\rm  \ldots\ Instructions  (indentées) exécutées si les expressions}\par 
\vspace{\neglskip}\mbox{\blc\ \ \ \ \ \ \ \small\rm  booléennes 1, \ldots, $n$ valent {\mcl False} \ldots}\par 
\medskip}
\end{minipage}\\ \hline 
\end{tabular}
\end{center}

 Le mot {\mcl elif} est une contraction de {\mcl else}~{\mcl if} (sinon si). Les expressions
 booléennes sont évaluées dans l'ordre et seules les instructions de la première alternative
 dont l'expression vaut {\mcl True} sont exécutées.
 Cette forme de test multiple est illustrée ci-dessous dans le programme construisant la frise suivante~:

  {\raggedright 
 \includevspacePostScript[1.0cm]{FriseDeVaguelettes}{0.95}{0.95}
 \input python/FriseDeVaguelettes.tex
 }

Pour terminer nous allons introduire de nouvelles commandes du \robot\ qui permettent
de faire pivoter le \robot\ sur place vers une orientation donnée. Par exemple la commande
{\idt pn()} qui permet de faire \underline{p}ivoter le \robot\ vers le \underline{n}ord peut
être définie à l'aide de la procédure \python\ suivante~:

{
\parindent \parindentprog
\medskip

\vspace{\neglskip}\mbox{\blc {\cmt \#\ Faire\ pivoter\ le\ robot\ sur\ place\ vers\ le\ nord}}\par
\vspace{\neglskip}\mbox{\blc {\mcl def}{\idt \ pn():}}\par
\vspace{\neglskip}\mbox{\blc \blk\blk{\mcl if}{\idt \ tne():}}\par
\vspace{\neglskip}\mbox{\blc \blk\blk\blk\blk{\idt phtg()}}\par
\vspace{\neglskip}\mbox{\blc \blk\blk{\mcl elif}{\idt \ te():}}\par
\vspace{\neglskip}\mbox{\blc \blk\blk\blk\blk{\idt pqtg()}}\par
\vspace{\neglskip}\mbox{\blc \blk\blk{\mcl elif}{\idt \ tse():}}\par
\vspace{\neglskip}\mbox{\blc \blk\blk\blk\blk{\idt p3htg()}}\par
\vspace{\neglskip}\mbox{\blc \blk\blk{\mcl elif}{\idt \ ts():}}\par
\vspace{\neglskip}\mbox{\blc \blk\blk\blk\blk{\idt pdt()}}\par
\vspace{\neglskip}\mbox{\blc \blk\blk{\mcl elif}{\idt \ tsw():}}\par
\vspace{\neglskip}\mbox{\blc \blk\blk\blk\blk{\idt p3htd()}}\par
\vspace{\neglskip}\mbox{\blc \blk\blk{\mcl elif}{\idt \ tw():}}\par
\vspace{\neglskip}\mbox{\blc \blk\blk\blk\blk{\idt pqtd()}}\par
\vspace{\neglskip}\mbox{\blc \blk\blk{\mcl elif}{\idt \ tnw():}}\par
\vspace{\neglskip}\mbox{\blc \blk\blk\blk\blk{\idt phtd()}}\par

\medskip }

\noindent Le programme {\idt OrientationAuNord.py} du dossier {\idt livre/\allowbreak 15\_if} essaie cette
procédure dans les huit directions. Nous pourrions définir de la même façon les  commandes de pivotement du \robot\
sur place dans l'une des huit directions principales de la rose des vents qui sont données dans la table
\ref{tab:pnetc} ci-dessous.

\begin{table}[htbp]
\begin{center}
\begin{tabular}{@{}|l|p{11.7cm}|@{}} \hline
\multicolumn{2}{|l|}{\em Commandes de pivotement du \robot\ sur place~:} \\ \hline \hline
{\idt pn()} & fait \underline{p}ivoter le \robot\ sur place vers le nord (\underline{n}). \\ \hline
{\idt pne()} &  fait \underline{p}ivoter le \robot\ sur place  vers le nord-est (\underline{ne}). \\ \hline
{\idt pe()} &  fait \underline{p}ivoter le \robot\ sur place vers l'est (\underline{e}). \\ \hline
{\idt pse()} &  fait \underline{p}ivoter le \robot\ sur place vers le sud-est (\underline{se}). \\ \hline
{\idt ps()} &  fait \underline{p}ivoter le \robot\ sur place vers le sud (\underline{s}). \\ \hline
{\idt psw()} &  fait \underline{p}ivoter le \robot\ sur place vers le sud-ouest (\underline{sw}). \\ \hline
{\idt pw()} &  fait \underline{p}ivoter le \robot\ sur place vers l' ouest (\underline{w}). \\ \hline
{\idt pnw()} &  fait \underline{p}ivoter le \robot\ sur place vers le nord-ouest (\underline{nw}). \\ \hline
\end{tabular}
\end{center}
\caption{Pivotement du \protect\robot\ sur place dans une direction donnée\label{tab:pnetc}}
\end{table}

\noindent \rule{137mm}{0.1mm} %---------------------------------------------
\begin{exercice}{Écrire des programmes \python\  comportant des instructions
de test ``if'' pour réaliser les frises de la figure \ref{fig:frisif} ci-dessous.
\noindent $\Box$}\end{exercice}
\begin{figure}[htbp]
 \includevspacePostScript[3.8cm]{FriseIf}{0.95}{0.95}
 \caption{Frises de cercles et (semi-) octogones}\label{fig:frisif}
\end{figure}
\begin{corrige}{
Nous donnons le programme pour dessiner la frise d'octogones de la figure \ref{fig:frisif}
 (page \pageref{fig:frisif}), les autres se trouvant dans le dossier {\idt livre/\allowbreak 15\_if}~:
\input python/FriseDoctogones.tex
\noindent $\Box$}\end{corrige}
\vfill\pagebreak[3]

%%%%%%%%%%%%%%%%%%%%%%%%%%%%%%%%%%%%%%%%%%%%%%%%%%%%%%%%%%%%%%%%%%%%%%%%%%%
\noindent \framebox[\textwidth]{
\begin{minipage}{\myparboxwidth}
\setlength{\parindent}{2em}
\small \sf
\vspace{\abovedisplayskip}
\section*{Addenda}
\index{Test}
L'instruction de test {\idt si} \ldots {\idt alors} \ldots {\idt sinon} \ldots fut introduite vers la fin des années
1950 dans le langage de programmation {\sc algol 60}. Ce langage, conçu par un groupe
d'experts internationaux, eut une très grande influence dans la conception de langages de programmation
ultérieurs notamment {\sc algol w}, \pascal, {\sc ada},  {\sc modula} et bien d'autres.  
Une autre forme d'instruction de test par cas fut introduite dans le langage de programmation 
{\sc algol w} (également conçu par N. Wirth) et reprise en \pascal\ sous le nom de ``case''.
En \ocaml, l'instruction de test par cas s'appelle le {\em filtrage} (``match'' en anglais).
\python\ possède aussi, depuis sa version 3.10, une instruction {\idt match}, mais nous ne l'utiliserons
pas~: en \python, on écrit habituellement un test par cas avec une suite de {\mcl if} \ldots\ {\mcl elif} \ldots\
{\mcl else}, en comparant la valeur d'une expression à chacun des cas possibles.

Par exemple {\idt valOrientation()} est
égal au numéro de l'orientation du \robot\ (0 pour le nord, 1 pour le nord-est, 2 pour l'est, 
etc.). La bibliothèque {\idt LibRobot} déclare les constantes {\idt nord}, {\idt nordEst}, {\idt est},
{\idt sudEst}, {\idt sud}, {\idt sudOuest}, {\idt ouest} et {\idt nordOuest}, égales à $0$, $1$, \ldots, $7$.
On pourra écrire~:
 

\index{Test}
{
\parindent \parindentprog
\medskip

\vspace{\neglskip}\mbox{\blc {\mcl if}{\idt \ valOrientation()\ ==\ nord:\ \ \ \ \ \ \ \ \ \ }{\cmt \#\ nord}}\par
\vspace{\neglskip}\mbox{\blc \blk\blk{\idt ...}}\par
\vspace{\neglskip}\mbox{\blc {\mcl elif}{\idt \ valOrientation()\ ==\ nordEst:\ \ \ \ \ }{\cmt \#\ nord{-}est}}\par
\vspace{\neglskip}\mbox{\blc \blk\blk{\idt ...}}\par
\vspace{\neglskip}\mbox{\blc {\mcl elif}{\idt \ valOrientation()\ ==\ est:\ \ \ \ \ \ \ \ \ }{\cmt \#\ est}}\par
\vspace{\neglskip}\mbox{\blc \blk\blk{\idt ...}}\par
\vspace{\neglskip}\mbox{\blc {\mcl elif}{\idt \ valOrientation()\ ==\ sudEst:\ \ \ \ \ \ }{\cmt \#\ sud{-}est}}\par
\vspace{\neglskip}\mbox{\blc \blk\blk{\idt ...}}\par
\vspace{\neglskip}\mbox{\blc {\mcl elif}{\idt \ valOrientation()\ ==\ sud:\ \ \ \ \ \ \ \ \ }{\cmt \#\ sud}}\par
\vspace{\neglskip}\mbox{\blc \blk\blk{\idt ...}}\par
\vspace{\neglskip}\mbox{\blc {\mcl elif}{\idt \ valOrientation()\ ==\ sudOuest:\ \ \ \ }{\cmt \#\ sud{-}ouest}}\par
\vspace{\neglskip}\mbox{\blc \blk\blk{\idt ...}}\par
\vspace{\neglskip}\mbox{\blc {\mcl elif}{\idt \ valOrientation()\ ==\ ouest:\ \ \ \ \ \ \ }{\cmt \#\ ouest}}\par
\vspace{\neglskip}\mbox{\blc \blk\blk{\idt ...}}\par
\vspace{\neglskip}\mbox{\blc {\mcl elif}{\idt \ valOrientation()\ ==\ nordOuest:\ \ \ }{\cmt \#\ nord{-}ouest}}\par
\vspace{\neglskip}\mbox{\blc \blk\blk{\idt ...}}\par
\vspace{\neglskip}\mbox{\blc {\mcl else}{\idt :\ \ \ \ \ \ \ \ \ \ \ \ \ \ \ \ \ \ \ \ \ \ \ \ \ \ \ \ \ \ \ \ \ }{\cmt \#\ impossible}}\par
\vspace{\neglskip}\mbox{\blc \blk\blk{\mcl pass}}\par

\medskip }

\vspace{\abovedisplayskip}
\end{minipage}
}
%%%%%%%%%%%%%%%%%%%%%%%%%%%%%%%%%%%%%%%%%%%%%%%%%%%%%%%%%%%%%%%%%%%%%%%%%%%%
%%%%%%%%%%%%%%%%%%%%%%%%%%%%%%%%%%%%%%%%%%%%%%%%%%%%%%%%%%%%%%%%%%%%%%%%%%%
\noindent \framebox[\textwidth]{
\begin{minipage}{\myparboxwidth}
\setlength{\parindent}{2em}
\small \sf
\vspace{\abovedisplayskip}
\noindent On notera que la comparaison s'écrit avec deux signes égal `{\idt ==}', le signe
`{\idt =}' seul servant à déclarer des constantes. Avec
`{\mcl if}\ {\idt valOrientation()\ =\ nord:}', \python\ répondrait
`{\tt SyntaxError: cannot assign to function call here. Maybe you meant '==' instead of '='?}'
(on ne peut pas donner de valeur à un appel de fonction, vouliez-vous dire `{\idt ==}'$\:$?).
On peut aussi comparer deux nombres avec {\idt !=} (différent de, $\neq$), {\idt \char60{}}, {\idt \char60{}=} ($\leq$),
{\idt \char62{}} et {\idt \char62{}=} ($\geq$).
Le dernier cas {\mcl else} regroupe toutes les autres valeurs, comme le mot
{\mcl otherwise} (`autrement' en anglais) de certaines extensions de \pascal. 
L'instruction {\mcl pass} (`passer' en anglais) ne fait rien~; on peut aussi omettre ce dernier cas.
Plusieurs cas semblables peuvent être regroupés 
comme dans l'exemple ci-dessous, où l'on teste si le reste de la division du numéro d'orientation par 2 est nul~:


% exemple de regroupement de cas
{
\parindent \parindentprog
\medskip

\vspace{\neglskip}\mbox{\blc {\mcl if}{\idt \ valOrientation()\ \%\ 2\ ==\ 0:\ \ }{\cmt \#\ N,\ E,\ S,\ W}}\par
\vspace{\neglskip}\mbox{\blc \blk\blk{\idt ...}}\par
\vspace{\neglskip}\mbox{\blc {\mcl else}{\idt :\ \ \ \ \ \ \ \ \ \ \ \ \ \ \ \ \ \ \ \ \ \ \ \ \ \ }{\cmt \#\ autres\ cas}}\par
\vspace{\neglskip}\mbox{\blc \blk\blk{\idt ...}}\par

\medskip }

On notera qu'avec des tests {\mcl if}, les cas peuvent être des noms de constantes comme {\idt nord}
aussi bien que des nombres comme {\idt 0}, ce qui rend le programme plus lisible.
Ce n'était pas le cas avec le filtrage d'\ocaml\ où `{\idt \char124{}\ nord\ {-}\char62{}}' ne signifie pas
`si l'orientation est le nord' mais déclare un nouvel identificateur {\idt nord} qui convient pour
toutes les valeurs (il en est de même avec le {\idt match} de \python).
 
\vspace{\abovedisplayskip}
\end{minipage}
}
%%%%%%%%%%%%%%%%%%%%%%%%%%%%%%%%%%%%%%%%%%%%%%%%%%%%%%%%%%%%%%%%%%%%%%%%%%%% 
\newpage                 % Si on termine le chapitre par une page blanche
\thispagestyle{empty}
